\section{From perceived risk to action: ecological and laboratory comparisons}
\label{sec:action}

A marmot that hears an alarm call can look up, stop feeding, or run for its burrow. Recognizing that a sound signals risk does not settle which of these it will do. Freezing, flight, vigilance, delayed re-emergence, and reduced foraging differ in cost and consequence, and laboratory and field settings differ in whether those costs exist at all. Comparing laboratory and field receivers shows how much the choice of action depends on those costs, and the defensive circuits of \secref{sec:mechanisms} offer hypotheses about the patterns seen in the wild (representative contrasts are in \tabref{tab:comparison}).

\subsection{Graded versus categorical responses}
\label{sec:action-graded}

Some rodents respond to alarm calls along a sliding scale; others choose among distinct actions. Yellow-bellied marmots treat call rate as information about risk, and adding harsh, noisy components to calls further suppresses foraging \citep{BlumsteinArmitage1997,BlumsteinRecapet2009}. Richardson's ground squirrels adjust vigilance differently to chirps, whistles, and chucks \citep{SloanEtAl2005}. Gunnison's prairie dogs are different: playback of alarm calls previously elicited by humans, dogs, or coyotes reproduces the matching escape patterns, whereas a control matched in rhythm, frequency, and amplitude does not \citep{KiriazisSlobodchikoff2006}. These are different informational solutions. Treating all sciurid alarms as predator labels, or all as urgency signals, misreads the comparative evidence (\secref{sec:features}).

Predatory imminence theory offers one way to relate the two patterns. It treats vigilance, freezing, and escape as distinct defensive modes that an animal enters as a threat is judged nearer, rather than as points on a single intensity scale \citep{Fanselow1994}. On that view, graded vigilance adjusts behavior within a mode, whereas categorical escape selects between modes. The framework is functional; whether sciurids implement it with the midbrain circuits studied in laboratory mice is unknown.

\subsection{Action choice depends on context and state}
\label{sec:action-context}

The same sound can produce opposite actions. Mice exposed to trains of 17--20~kHz frequency sweeps at 85~dB either flee or freeze; the two responses are negatively correlated across animals, and cage familiarity and prior foot shock shift the balance between them \citep{MongeauEtAl2003}. The sound is synthetic, which is the point: the result isolates state and opportunity from any meaning carried by the call, and the trade-off between flight and freezing is what a competition between defensive output systems would produce (\secref{sec:mechanisms-amygdala}). Experience shows a similar dissociation for natural calls. Rats freeze selectively to 22-kHz calls only after aversive experience \citep{KimEtAl2010,EndresEtAl2007}, yet they slow their movement and heart without that history \citep{OlszynskiEtAl2020}.

Field receivers weigh refuge, food, and social setting in the same way. Alarm playback makes chipmunks more vigilant, shortens their runs to cover, delays re-emergence, and reduces intake \citep{BaackSwitzer2000}; gray squirrels hearing predator calls delay and reduce foraging without changing how far they stray from refuge \citep{JayneEtAl2015}.
% VERIFY: BaackSwitzer2000, JayneEtAl2015 (not full-text checked).
In both laboratory and field, prior pairing of sound with danger shapes later action: experience gates freezing to 22-kHz calls \citep{KimEtAl2010}, and a receiver's history with particular callers revises its vigilance \citep{HareAtkins2001,BlumsteinEtAl2004}. The question is shared; the circuits need not be.

\subsection{What defensive circuits suggest about field responses}
\label{sec:action-circuits}

To our knowledge, no study has causally established the circuits that mediate natural alarm-call responses in sciurids. Playback has revealed region-specific neural activation in captive ground squirrels \citep{FreemanEtAl2020}, but recordings or manipulations of defensive circuits during alarm responses in free-ranging sciurids remain unreported. Such recordings are becoming feasible: an implantable platform that logs brain activity and movement has been demonstrated in a captive fox squirrel \citep{GaidicaDantzer2022}. The laboratory circuitry is still useful for field biology, because it shows that different behavioral outcomes can place different demands on the receiver, and those demands make predictions that field experiments can test. Three contrasts stand out. Each is a hypothesis; none has been shown to rest on the same circuits in sciurids.

\textbf{Harsh, abrupt sounds and graded vigilance.} In laboratory mice, loud or looming sounds trigger defense through fast, largely experience-independent routes---the brainstem startle pathway and cortical projections to the IC and SC---without first assigning learned meaning \citep{DavisEtAl1982,XiongEtAl2015,LiEtAl2021}. Marmots respond more strongly to calls with added noisy, nonlinear components \citep{BlumsteinRecapet2009}. If those components act through a similar arousal route, their effect should appear in inexperienced animals, generalize across callers, and decline like sensory habituation. If instead marmots have learned that harsh calls signal higher risk, the effect should depend on experience and be specific to callers whose calls have been paired with danger. Rate-based urgency information \citep{BlumsteinArmitage1997} could follow either pattern; separating harshness from rate would distinguish them.

\textbf{Predator-specific escape and action selection.} Prairie-dog escape that matches the predator type \citep{KiriazisSlobodchikoff2006} requires the receiver to evaluate a call category and select among distinct motor programs. Laboratory circuits show how such selection could work: competing CeA populations bias conditioned freezing versus flight \citep{FadokEtAl2017}, and separate cortical outputs to the striatum and the SC support freezing and flight to the same looming sound \citep{LiEtAl2021}. These results suggest a prediction for the field. If call evaluation and action selection are separable, reducing access to refuge should change the escape topography a given call elicits while leaving signs of recognition---orientation, vigilance, or cardiac change---intact.

\textbf{Vigilance, freezing, and the heart.} Vigilance (an upright scanning posture that trades feeding time for information) and freezing (immobility) are often scored as alternative ``fear'' responses. In laboratory animals, freezing comes with a parasympathetic slowing of the heart that may aid perception and preparation for action, whereas flight comes with sympathetic acceleration \citep{Roelofs2017}. Cardiac discrimination of alarm call types in juvenile ground squirrels \citep{Mateo1996} shows that heart rate can be measured in field-relevant receivers. Whether alarm-evoked vigilance carries the cardiac signature of freezing, of flight, or neither is unknown, and the answer would show whether vigilance is a form of freezing adapted for information gathering or a distinct mode.

\subsection{Heterospecific information and safety signals}
\label{sec:action-hetero}

Rodents listen to other species as well as their own. Marmots and ground squirrels respond to heterospecific alarms much as they do to conspecific alarms, relative to birdsong controls \citep{Shriner1998}; Cape ground squirrels are as vigilant to crowned-lapwing alarms as to their own species' alarms \citep{WatermanMai2020}; and gray squirrels respond more to familiar bird alarms, and to unfamiliar alarms that sound similar, than to unfamiliar dissimilar alarms \citep{GetschowEtAl2013}.
% VERIFY: Shriner1998, WatermanMai2020, LillyEtAl2019 (not full-text checked where marked).
Sounds can also signal safety. After a hawk call, ordinary bird chatter lets squirrels relax vigilance faster than ambient noise does \citep{LillyEtAl2019}. And predator source is a poor guide to impact: chipmunks change foraging more to titmouse mobbing calls than to hawk calls \citep{SchmidtEtAl2008}. Social ecology expands the set of meaningful sounds, but it also confounds familiarity with acoustic similarity unless both are manipulated.

\subsection{Costs, consequences, and where lab--field translation breaks down}
\label{sec:action-costs}

Vigilance, freezing, escape, slow recovery, and reduced foraging are not equivalent biological effects. Vigilance buys information at the cost of food; escape interrupts feeding and may expose the animal; a hormonal response without visible defense may still reflect assessment \citep{EilamEtAl1999}. A laboratory cage that offers neither food nor refuge cannot show ecological impact, however strong the freezing. Nor can ecological choices be read from neural activity in a conditioned laboratory animal, or from a single peripheral measure such as heart rate or c-Fos (\secref{sec:mechanisms}). In the other direction, field playback without matched acoustic controls or measured sound level at the receiver cannot attribute effects to particular acoustic features (\secref{sec:features}).

Temporal structure links the two settings in a similar way. Laboratory startle and amygdala responses habituate less to irregularly than to regularly timed sounds \citep{Davis1970,HerryEtAl2007}, and field responses recover when the caller changes \citep{Hare1998}. Both concern how the history of a sound shapes the next response (\secref{sec:persistence}), but laboratory timing effects with artificial sounds do not establish how natural calls should be timed to sustain field responses.

\subsection{Predator calls to prey, and the limits of attraction evidence}
\label{sec:action-scope}

Predator calls played to rodents give conflicting results in predator-naïve laboratory stock \citep{HendrieNeill1991,KindermannEtAl2009} and more consistent but still species-variable results in field-experienced animals \citep{HendrieEtAl1998,EilamEtAl1999}. Rat 50-kHz playback provides the appetitive contrast: approach and a transient rise in accumbens dopamine relative to matched noise \citep{WilluhnEtAl2014,WohrSchwarting2007}. That is attraction of laboratory rodents to conspecific calls, not evidence for field attraction in sciurids. The comparison here is how rodents respond to the sounds of predators and of conspecifics; how predators respond to the calls of their prey is a different problem.

% Main-text summary of recurring principles. Dense comparison: supplement/table_comparison_detail.tex.
\begin{table}[tbp]
\centering
\footnotesize
\singlespacing
\setlength{\tabcolsep}{5pt}
\renewcommand{\arraystretch}{1.3}
\caption{Recurring principles across laboratory and field rodent receivers. Each row states one point, the evidence that best illustrates it, and what that evidence does not establish. Fuller comparisons are in \stabref{tab:comparison-supp}{1}; study details are in \stabref{tab:evidence-supp}{2}.}
\label{tab:comparison}
\begin{tabular}{@{}>{\raggedright\arraybackslash}p{1.35in}>{\raggedright\arraybackslash}p{3.0in}>{\raggedright\arraybackslash}p{1.75in}@{}}
\toprule
\textbf{Principle} & \textbf{Illustrative evidence} & \textbf{What it does not show} \\
\midrule
Simple acoustic features often suffice
  & Tones approximating 50-kHz calls draw rats to the speaker, and frequency-matched tones reproduce most amygdala discrimination between call types \citep{WohrSchwarting2007,ParsanaEtAl2012a}.
  & That natural calls carry no further meaning. \\
Alarm calls carry different kinds of information
  & Marmot alarms grade risk by call rate; prairie-dog alarms specify escape suited to the predator \citep{BlumsteinArmitage1997,KiriazisSlobodchikoff2006}.
  & A single code shared across sciurids. \\
Experience gates action more than perception
  & Rats freeze to 22-kHz calls only after aversive experience, yet naïve rats still slow their heart and movement \citep{KimEtAl2010,OlszynskiEtAl2020}.
  & Whether predator calls are recognized without experience; tests conflict. \\
The body can respond without behavior
  & Spiny mice raise cortisol to owl calls without freezing or flight; young ground squirrels discriminate alarms by heart rate before they do so behaviorally \citep{EilamEtAl1999,Mateo1996}.
  & Which pathway produces the physiological change. \\
Parallel brain pathways allow these dissociations
  & Separate circuits support learned association, freezing, flight, and blood-pressure change \citep{LeDouxEtAl1988,FadokEtAl2017,TovoteEtAl2016,LiEtAl2021}.
  & That natural alarms use these circuits; most were mapped with tones or synthetic sounds. \\
Context selects the action
  & The same sound evokes flight or freezing depending on setting, and field receivers trade vigilance against foraging and refuge \citep{MongeauEtAl2003,BaackSwitzer2000}.
  & How costs are weighed in the brain. \\
Receivers track caller reliability
  & Ground squirrels discount unreliable callers, whereas marmots scrutinize them more closely \citep{HareAtkins2001,BlumsteinEtAl2004}.
  & Why the direction differs between species. \\
Declining responses have several causes
  & Responses recover when the caller changes, and irregular timing slows habituation to artificial sounds \citep{Hare1998,Davis1970,HerryEtAl2007}.
  & Which process underlies a given decline curve. \\
Meaningful sounds need not be threatening
  & Rats approach 50-kHz calls with a rise in accumbens dopamine, and oxytocin in auditory cortex enables mice to retrieve calling pups \citep{WilluhnEtAl2014,MarlinEtAl2015}.
  & Attraction to calls in free-living sciurids. \\
\bottomrule
\end{tabular}
\end{table}

